%------------------------------------------------------------------------------%
\section{Números Complexos}
\label{sec:NumerosComplexos}

Essa seção introduz o conceito de \emph{Números Complexos}\index{Número complexo}
e apresenta os resultados utilizados nesse material.

%Uma apresentação mais completa dos Complexos pode ser
%encontrada em [\emph{CITAR ALGUÉM}].

Como uma motivação para a criação dos números complexos considere a equação
\[
  x^2 + 1 = 0
\]
manipulando formalmente obtemos a solução
\[
  x = \sqrt{-1}
\]
Porém, para qualquer número real $a$ sabemos que \(a^2 \geq 0\) o que implica que não existe
o número \(\sqrt{-1}\) dentro do conjunto dos números reais e portanto a equação não possui solução.

Uma forma para contornar essa limitação é criar um novo número, denotado por $i$ e que chamamos de
\emph{Unidade imaginária}\index{Unidade Imaginária}, de forma que
\[
  i^2 = -1
\]
Além disso, determinamos que as operações com esse número obedecem as mesmas regras das operações com
os números reais. Por exemplo, podemos calcular as potências de $i$
\begin{align*}
	i^2 & = -1                    \\
	i^3 & = i^2i = (-1)i = -i     \\
	i^4 & = i^2i^2 = (-1)(-1) = 1 \\
	i^5 & = i^4i = i              \\
	i^6 & = i^5i = i^2 = -1
\end{align*}

Como o número $i$ não pertence aos reais estamos criando um novo conjunto numérico conhecido como
o \emph{Conjunto dos Números Complexos}\index{Número complexo!conjunto}
\[
  \C = \big\{\, a+bi \;\big|\; a, b \in \R \,\big\}
\]
Note que, identificarmos os números reais, $a$, com os números complexos \((a+0i)\) e
dessa forma entendemos os números reais como um subconjunto dos complexos.

No números complexos as operações de soma e produto são definidas
de modo compatível com as operações nos números reais.
A \emph{soma}\index{Número complexo!soma} de dois números complexos $v=a+bi$ e $w=c+di$ é definida por
\[
  v+w = (a+bi) + (c+di) = (a+c) + (b+d) i
\]
Enquanto o \emph{produto}\index{Número complexo!produto} é definido por
\begin{align*}
	vw & = (a+bi) \cdot (c+di)             \\
	   & = a \cdot (c+di) + bi \cdot(c+di) \\
	   & = ac + adi + bci + bdi^2          \\
	   & = ac + adi + bci - bd             \\
	   & = (ac-bd) + (ad+bc) i
\end{align*}
Propriedades das operações de números complexos. Sejam $u$, $v$ e $w$ números complexos, então
\begin{enumerate}[label=\alph*)]
\item \( v + w = w + v \)
\item \( v + (w + u) = ( v + w ) + u \)
\item \( v + 0 = 0 + v = v \) sendo que \( 0 = 0 + 0i \)
\item Para todo \(v=a+bi \in\C\), o número
      \[-v=-a-bi \qquad \text{é tal que} \qquad v+(-v) = 0 \]
\item \(vw = wv\)
\item \( v(wu) = (vw)u\)
\item \( 1v = v\) sendo que \(1=1+0i\)
\item Se \( v = a+bi \neq 0 \) então
      \[ v^{-1} = \frac{a}{a^2+b^2} - \frac{b}{a^2+b^2}\,i
        \qquad\text{é tal que} \qquad
        vv^{-1} = v^{-1}v = 1
      \]
\item \( v(w+u) = vw + vu  \)
\end{enumerate}
Como as operações de soma e produto nos complexos possui a essas propriedades
dizemos que o conjunto dos números complexos, com essas operações, é um \emph{corpo}\index{Corpo}.
Os números reais, $\R$, e os racionais, $\Q$, também são corpos.

\begin{example}
  Calculo das raízes da equação quadrática.
  \solution
  Considere a seguinte equação algébrica
  \[ ax^2 + bx + c = 0 \]
  onde $a\neq 0$. Sabemos que suas soluções são dadas pela fórmula de Bhaskara
  \[ x = \frac{-b \pm \sqrt{\Delta}}{2a} \]
  onde
  \[ \Delta = b^2 - 4 ac \]

  Quando os coeficientes $a$, $b$ e $c$ são reais
  a função \(y(x) = ax^2 + bx + c \)
  representa uma parábola do plano e as soluções da equação quadrática são as os pontos onde o gráfico
  cruza o eixo $x$. Nesse contexto, buscamos soluções reais e quando \(\Delta<0\) dizemos que o
  gráfico não cruza o eixo $x$.

  Porém, usando os números complexos, podemos calcular a raiz quadrada de números negativos obtendo
  \[ \sqrt{\Delta} = \sqrt{\abs{\Delta}}\;i \]
  e as soluções para a equação quadrática são o par conjugado
  \[ x_1 = -\frac{b}{2a} - \frac{\sqrt{\abs{\Delta}}}{2a}\,i \]
  \[ x_2 = -\frac{b}{2a} + \frac{\sqrt{\abs{\Delta}}}{2a}\,i \]

  É possível mostrar que no, conjunto dos números complexos,
  todo polinômio de grau $n$ possui exatamente $n$ raízes.
  Além disso, se os coeficientes do polinômio forem números reais
  então suas raízes serão números reais ou números complexos conjugadas.
\end{example}

Assim como representamos os números reais em uma reta, a Reta Real, representamos
os complexos no \emph{Plano Complexo}\index{Plano complexo}. Cada número $u=x+yi$ é associado ao ponto
$(x, y)$ do plano. A Figura~\ref{fig:PlanoComplexo} ilustra essa associação.

%------------------------------------------------------------------------------%
\begin{figure}
  \centering
  % Title: gl2ps_renderer figure
% Creator: GL2PS 1.4.2, (C) 1999-2020 C. Geuzaine
% For: Octave
% CreationDate: Fri Sep 10 20:09:38 2021
\setlength{\unitlength}{1pt}
\begin{picture}(0,0)
\includegraphics[width=14cm,height=5.2cm]{fig_complex_plane}
\end{picture}%
\begin{picture}(396,147)(0,0)
\fontsize{12}{0}\selectfont\put(274.961,63.0945){\makebox(0,0)[l]{\textcolor[rgb]{0,0,0}{{$x$}}}}
\fontsize{12}{0}\selectfont\put(81.0709,129.425){\makebox(0,0)[r]{\textcolor[rgb]{0,0,0}{{$y$}}}}
\fontsize{12}{0}\selectfont\put(81.0709,17.1732){\makebox(0,0)[r]{\textcolor[rgb]{0,0,0}{{$-y$}}}}
\fontsize{12}{0}\selectfont\put(280.063,132.997){\makebox(0,0)[l]{\textcolor[rgb]{0,0,0}{{$z=x+iy$}}}}
\fontsize{12}{0}\selectfont\put(280.063,13.6016){\makebox(0,0)[l]{\textcolor[rgb]{0,0,0}{{$\bar{z}=x-iy$}}}}
\fontsize{12}{0}\selectfont\put(142.299,82.4835){\makebox(0,0)[l]{\textcolor[rgb]{0,0,0}{{$\varphi$}}}}
\fontsize{12}{0}\selectfont\put(142.299,64.115){\makebox(0,0)[l]{\textcolor[rgb]{0,0,0}{{$\varphi$}}}}
\fontsize{12}{0}\selectfont\put(188.22,114.118){\makebox(0,0){\textcolor[rgb]{0,0,0}{{$r$}}}}
\fontsize{12}{0}\selectfont\put(188.22,32.4803){\makebox(0,0){\textcolor[rgb]{0,0,0}{{$r$}}}}
\end{picture}

  \caption{Representação dos Números Complexos em um Plano}
  \label{fig:PlanoComplexo}
\end{figure}
%------------------------------------------------------------------------------%

Além das operações de soma e produto, também definimos outras operações
para os complexos.

\begin{definicao}{Operações com Números Complexos}{def:complex-operations}
  Para um número \( z = x+yi \in \C \) definimos as operações:
  \begin{description}
  \item[Parte real]       \tabto{35mm} \( \Re(z) = x \)\index{Número complexo!parte real}
  \item[Parte imaginária] \tabto{35mm} \( \Im(z) = y \)\index{Número complexo!parte imaginária}
  \item[Conjugado]        \tabto{35mm} \( \overline{z} = x - yi \)\index{Número complexo!conjugado}
  \item[Módulo]           \tabto{35mm} \( \abs{z} = r = \sqrt{x^2+y^2\,} \)\index{Número complexo!módulo}
  \item[Argumento]        \tabto{35mm} \( \arg{z} = \varphi = \arctan\frac{y}{\,x\,}\)\index{Número complexo!argumento}
  \end{description}
\end{definicao}
A Figura~\ref{fig:PlanoComplexo} mostra a interpretação geométrica dessas operações.

É possível demonstrar as seguintes propriedades desses operações, para os números
complexos $v$ e $w$.
\begin{enumerate}[label=\alph*)]
\item \( v\overline{v}  = \abs{v}^2 \)
\item \( \overline{v+w} = \overline{v} + \overline{w} \)
\item \( \overline{vw}  = \overline{v}\overline{w} \)
\end{enumerate}

Um estudo do cálculo em variáveis complexas está além do escopo desse texto,
porém, para podemos introduzir a Transformada de Fourier, vamos precisar da definição de algumas
funções com variáveis complexas.

As funções algébricas, aquelas que utilizam apenas operações como soma, produto, e potências,
não precisam de uma definição especial, basta realizar as operações definidas para os complexos.
Porém, para as demais funções precisamos de outras técnicas. Para as funções analíticas,
isso é, as funções cujas Séries de Taylor convergem para a própria função, podemos usar essa série
para definir as funções nos complexos. As funções que vamos utilizar com maior frequência
são o seno, cosseno e a exponencial
\begin{align}
	\sin z & = z - \frac{z^3}{3!} + \frac{z^5}{5!} + \frac{z^7}{7!} - \cdots      \\[3mm]
	\cos z & = 1 - \frac{z^2}{2!} + \frac{z^4}{4!} + \frac{z^8}{8!} - \cdots      \\[3mm]
	e^z    & = 1 + z + \frac{z^2}{2!} + \frac{z^3}{3!} +  \cdots
\end{align}

Comparando as séries de Taylor dessas funções podemos demonstrar a \emph{Identidade de Euler}\index{Identidade de Euler}
que associa a exponencial com as funções seno e cosseno
\begin{equation}
  e^{i\theta} = \cos(\theta) + i\sin(\theta)  \qquad \forall \; \theta \in\R
\end{equation}
Com essa identidade podemos interpretar a \emph{Exponencial Complexa}\index{Exponencial Complexa}
como sendo uma oscilação com decrescimento ou decaimento exponencial, pois
\[e^{z} = e^{a+bi} = e^a e^{bi} =e^a  \big[ \cos(b) + i\sin(b) \big] \]
Dessa forma \(e^a\) determina o crescimento ou decrescimento, enquanto
\(\cos(b) + i\sin(b)\) determina a oscilação.

%------------------------------------------------------------------------------%
\subsection*{Cálculo com Funções e Variáveis Complexas}

Ao longo desse texto diversas vezes utilizamos funções de uma variável real
que assumem valores complexos, $f:D\subset\R\to\C$, nesses casos podemos
aplicar as técnicas do Cálculo de funções reais considerando a $i$ como uma
constante, por exemplo,
\begin{gather*}
  \frac{d}{dx}\left(ax^2 + ix + 3i\right) = 2ax + i                    \\[3mm]
  \frac{d}{dx}\sin(ix) = i\cos(ix)                                     \\[5mm]
  \int ax^2 + ix + 3i \,dx = \frac{ax^3}{3} + \frac{ix^2}{2} + 3ix + C \\[3mm]
  \int \sin(ix) \,dx = -\frac{\cos(ix)}{i} = i\cos(ix)
\end{gather*}

Porém, \emph{não} podemos utilizar as mesmas regras para funções
de variáveis complexas, $f:D\in\C\to\C$. Para esses casos precisamos da
teoria própria para as variáveis complexas.

%------------------------------------------------------------------------------%
